\section{Separate finite-class observation-selection study}
\label{app:observation-selection}

\subsection{Development history and a changed estimand}

Earlier prediction criteria remain failed. In one 1,600-draw study,
candidate consensus issued 639 point predictions, 189 incorrect
(29.58\% conditional error). A history filter chosen from those data
and tested on a separate 1,600-draw sample issued 436 predictions with
94 errors (21.56\%). Both failed the joint target of at most 5\%
conditional error and at least 25\% prediction coverage. The present
certificate comparison has a different estimand.

Reanalysis of the second sample with a covering class found 31 initially
singleton targets among 1,600 records. A response-branch scan on those
records found 134 possible gains and zero possible losses for the
selected policy versus the fixed query, but collected no observations.
It did not establish realized superiority. The policy was chosen with
these results known, then frozen before the new class-frame sample.
Development records are excluded from the new denominator.

\subsection{Class, target and sampling contract}

The candidate class contains deterministic systems with at most three
states and binary inputs and outputs, covered by 46,656 labeled
three-state tables. Labeled candidates and sampled behavior classes
are different units. A certificate requires a nonempty compatible set
unanimous about the specified two-output target from the original
cutoff. Coverage and implementation correctness are premises, not
established for arbitrary worlds by observing no mistakes. Actors receive
the legal history, target definition and candidate class, without true
system tables or state. The observation tool returns only a permitted
intermediate response, never the target continuation.

Write the transition/output table as $(\delta(s,a),\lambda(s,a))$,
with $s\in\{0,1,2\}$, $a\in\{0,1\}$ and initial state 0. The table
emits $\lambda(s,a)$ on a transition from $s$ and moves to
$\delta(s,a)$. The eight observed input--output pairs determine each
compatible candidate's cutoff state $s_8$. Query $q_b$ applies just
input $b$ to a fresh copy of that cutoff and returns
$\lambda(s_8,b)$, for $b\in\{0,1\}$. The target is
\[
T(s_8)=\bigl(\lambda(\delta(s_8,0),0),\,
                 \lambda(\delta(s_8,1),0)\bigr).
\]
These are the second outputs of $(0,0)$ and $(1,0)$, respectively,
starting from separate copies of the original cutoff. The first outputs
are not the target; the chosen query does not shift its reference state.
The stored target code is $2T_0+T_1\in\{0,1,2,3\}$. The fixed arm
always chooses $q_0$; $q_0$ also wins the selected policy's final tie.
Only $q_0$ and $q_1$ are legal: a query cannot request the second step
or the target pair. Actor-visible definitions and candidate-side
computation are disclosed in \texttt{artifact/inspection/wire.py} and
\texttt{compute.py}; they do not supply the true system's table.

Canonicalization gives 11,500 behaviors in the full at-most-three-state
class: 4 with one minimal reachable state, 144 with two, and 11,352 with
three. The original sampling design restricts worlds to exactly three
minimal reachable states, removing the 148 lower-state behaviors. Seven
earlier development fixtures represent two distinct eligible three-state
behaviors; removing these gives the frozen 11,350-class catalog. The first
1,600-draw development sample uses 1,499 distinct classes, leaving 9,851;
the second uses 1,479 additional classes, leaving 8,372. Each has 256
eight-input histories. Thus the apparent 150-class gap is an explicit
sampling restriction and prior-fixture exclusion, not a change to the
actor's covering hypothesis class. The two classes and all subsequent
exclusions are reconstructible by canonical-table identity and index in
\texttt{artifact/acquisition/replay.py}. The original local structural
intent is dated 26 October 2026, before the 27 October main freeze;
these local dates and hashes are not independent public time attestations.
The fixed sample contains 5,000 draws, uniform with replacement over
their 2,143,232 pairs. The implementation uses an operating-system seed,
a SHA-256 counter stream and rejection mapping to avoid index-modulo
bias. Six upper-residue rejections among 5,006 stream words are sampling
implementation events, not outcome exclusions. The inferential model
is independent uniform draws from this finite frame, not a claim that
cryptographic expansion provides physically independent randomness.
All seven duplicate-pair excess draws and 97 initially singleton targets
remain in $N$; 4,903 draws are initially ambiguous. Visible-history
overlap is allowed. Class exclusion does not imply absence from catalog
construction or all prior research.

\subsection{Pre-specified interval, size and decision}

Let $G$ and $L$ count selected-only and fixed-only certificates, and
$D=G+L$. Write $\theta$ for discordance probability and $\beta$ for the
selected-only probability conditional on discordance. The difference
in marginal certificate probabilities is
\[
\Delta=\theta(2\beta-1),\qquad
\widehat\Delta=(G-L)/N.
\]
The primary procedure forms central 97.5\% Clopper--Pearson intervals
for $\theta$ from $D/N$ and for $\beta$ from $G/D$ (full $[0,1]$ when
$D=0$), then takes the minimum and maximum of
$\theta(2\beta-1)$ at their four rectangle corners. Conditional
coverage for the second interval, unconditional coverage for the first
and the union bound give joint coverage of at least 95\% under the
paired-draw model. This is a conservative projection, not the same
procedure as a marginal Wald interval. The numerical implementation
and boundary cases were checked before sampling.

The pre-specified $N$ grid was 1,000 through 6,000 in steps of 1,000.
At planning difference 0.02 and discordance 0.15 (selected-only 0.085,
fixed-only 0.065), $N=5{,}000$ was the first to reach the 90\% power
target: 91.65415\%, versus 83.92156\% at 4,000. Planning used the same
zero-exclusion decision as the primary interval. For positive
discordance, it requires the central 97.5\% conditional interval for
$\beta$ to exclude one half, rather than substituting a 5\% two-sided
test. Power depends on the assumed discordance; 0.15 was not a proven
upper bound. Two percentage points was a research-interest threshold,
not a utility valuation or a claim that the true effect was that large.

One final readout followed all commitments and the target audit.
With $G=90$, $L=1$, the component intervals are
$\theta\in[0.0142214,0.0229103]$ and
$\beta\in[0.9319176,0.9998618]$. The projected interval is
$[0.01228495885,0.02290400123]$. Its lower endpoint exceeds zero,
so the positive-difference claim is supported; its inclusion of 0.02
leaves the interest threshold unresolved. Technical missingness would
retain the fixed denominator rather than trigger replacement draws.
Any incorrect certificate, identity violation or target leak is a
validity failure independent of the confidence interval. There was no
outcome-based sample extension or policy revision.

\subsection{Execution controls, diagnostics and work}

Separate actor processes and state copies implement the two arms.
The controls exercise invalid-query rejection, arm-order invariance
of normalized scientific outputs, masked-target invariance, rejection
of target-continuation access, pre-audit commitments, incorrect-certificate
detection and complete work accounting. Five thousand paired predictions
were committed before privileged target evaluation. The 10,000 actor
processes and keep-awake process exited. Pre-main checks used 77
synthetic extra observations and 81 actor processes; they are excluded
from the scientific sample. Same-host checks and clean-copy replay are
not independent laboratory validation.

Among initially ambiguous draws, the fixed and selected policies
certified 186/4,903 and 275/4,903 targets. Among the 686 draws with
different queries, the counts were 31/686 and 120/686. These are
descriptive conditional rates. The one fixed-only result means that
counting singleton-target branches does not maximize certification for
every actual response. Likewise, the 4,627 neither-certified pairs
do not prove that every permitted one-observation method would fail.

\begin{table}[!ht]
\centering
\small
\begin{tabular}{lrr}
\toprule
Measured operation & Fixed & Selected\\
\midrule
Environment extra observations & 5,000 & 5,000\\
Candidate--history checks & 233,280,000 & 233,280,000\\
Candidate history steps & 494,337,048 & 494,337,048\\
Query--candidate evaluations & 5,283,246 & 10,566,492\\
Selection-branch checks & 0 & 20,000\\
Target--candidate evaluations & 5,283,246 & 5,283,246\\
Target--candidate steps & 21,132,984 & 21,132,984\\
Certificate-branch lookups & 5,000 & 5,000\\
Reused query values within an arm & 0 & 5,283,246\\
Actor CPU seconds & 136.38 & 137.96\\
Actor wall seconds & 143.91 & 146.28\\
Actor-process wall seconds & 311.74 & 313.80\\
\bottomrule
\end{tabular}
\caption{Work for the separate 5,000-draw acquisition study.
Counters are operation-specific and may overlap; reused values are
not additional evaluations. No cross-arm or cross-record caching is
used. Timing does not identify an order-independent speed difference.}
\label{tab:observation-work}
\end{table}


The total main-run elapsed time was 640.28 seconds. In addition to
10,000 extra observations, environment accounting records 40,000
history-generation steps, 5,000 resets, 10,000 clones, zero history-replay
steps and 20,000 target-audit steps. The audit compares 655 issued
certificates to their targets. These operation-specific counters are
not an additive cost score. Timing is descriptive because the fixed arm
always ran first. Model/provider requests were zero; research-assistant
use was unmeasured.


\subsection{Raw-input reader reconstruction}
The bundle includes the original catalogs, exclusions, seed, rejection
receipts, frozen plan, commitments, events, histories and arm records,
plus collector/environment source as inspection-only text. Plain JSON
columns and string references retain every original field and reproduce
all 5,000 original record hashes. Retained process IDs and timings establish byte identity only.

Run \texttt{python3.12 -I -B -S artifact/acquisition/replay.py} from the
bundle root. It reconstructs the universe, ordered catalogs, exclusions
and full first development manifest; verifies the original hashes and
5,006 sampling receipts; regenerates histories, both policies' query
responses, candidate computations and certificates; and fixes predictions
before recomputing targets. Every scientific field and work counter is
compared before recomputing the primary interval. The archived
commitment-before-truth event order is also checked.

This standard-library replay collects no new sample, calls no model, and
executes neither participant programs nor the historical isolation collector.
Independent historical attestation remains outside its scope.
